% !TeX encoding = UTF-8 !TeX spellcheck = en_US
\documentclass[11pt, a4paper]{article}

% The title of the current document to be produced.
\newcommand{\doctitle}{Assignment \#2: 29 + 13 = 42 points total }
\newcommand{\assignments}{\bluetext{\textbf{Assignment:}}}
%
\setlength{\unitlength}{1in}
\renewcommand{\arraystretch}{2}
\setlength{\columnsep}{2cm}

\newcommand\dunderline[3][-1pt]{{%
  \sbox0{#3}%
  \ooalign{\copy0\cr\rule[\dimexpr#1-#2\relax]{\wd0}{#2}}}}

\usepackage[sfdefault]{atkinson}
\usepackage[T1]{fontenc}

\def\code#1{\bluetext{\texttt{#1}}}


%------------------------------------------------------------
% Import commands for both teacher and course information.  | NOTE: Change your
% teacher and course info in these files. |
%------>------>------>------>------>------>------>------>-->|
%
%------------------------------------------------------------
%-- Import packages and custom command definitons.          |
%------>------>------>------>------>------>------>------>-->|
%% ==================================
%% ===== Teacher info             ===
%% ==================================
%%
\newcommand{\instructor}{Katelyn Breivik}
\newcommand{\office}{Wean 8319}
\newcommand{\hours}{3pm on Mondays}
\newcommand{\phone}{}
\newcommand{\college}{}
\newcommand{\email}{kbreivik@andrew.cmu.edu}
\newcommand{\faculty}{}
\newcommand{\department}{Department of Physics}
                              %|
%% ==================================
%% ===== Course-specific commands ===
%% ==================================

%- Instructions: change course info here. 
\newcommand{\semester}{Fall 2023}

\newcommand{\ponderation}{2-4-3 (Theory-Lab-Homework)}
\newcommand{\coursetitle}{Astrophysics of Stars and the Galaxy}
\newcommand{\coursenumber}{[33-467]}
\newcommand{\prerequisite}{33-224, 33-234}
 
\input{../includes/packages-imports}                          %|  
\input{../includes/custom-commands}   
%
%---> Genereate & inject metadata                           |
\input{../includes/hyperef.doc.info}                          %|
%------------------------------------------------------------

\topmargin -70pt
\begin{document} 

\normalfont

%-------------------------------------------------------------
%-- Make the header of the document                          |
%------>------>------>------>------>------>------>------>--> |
%---------------------------------------------------------------------
%-- The following produces the document header including the title.  |
%---------------------------------------------------------------------

\noindent % <-- need to have this first.

\hfill	
\begin{minipage}{0.60\textwidth}%
	\raggedleft%
	{\Large \textsc{\doctitle}\par}
	\doublerule 
\end{minipage}%
\vspace{0.9cm}

  
%-------------------------------------------------------------
%-- Problem # 1                                              |
%------>------>------>------>------>------>------>------>--> |
\noindent \dunderline{1.0pt}{\textbf{Problem \#1: Stellar timescales (Based on
 Pols notes 2.5) 29 points total}}\\
\noindent By the end of this problem, you should have a solid understanding of
 the different timescales of how stars respond due to changes of their
 hydrostatic or thermal equilibrium as well as due to nuclear burning in their
 centers.

 \vspace{4pt}
 \noindent \textbf{Problem \#1a: The nuclear timescale ($\tau_{\rm{nuc}}$): 8
 points total}
 
 \vspace{4pt}
 \noindent ($i$) Calculate the total mass of hydrogen available for fusion over
 the lifetime of the Sun, if $70\%$ of its mass was hydrogen when the Sun was
 formed, and only $13\%$ of all hydrogen is in the layers where the temperature
 is high enough for fusion. \textbf{2 pts; 1 pt for trying; 1 pt for correct
 answer; -0.5 for units}

 \vspace{8pt}

 \noindent ($ii$) Derive an expression for the nuclear timescale in solar units,
 i.e. expressed in terms of $(R/R_{\odot}$), $(M/M_{\odot}$), and
 $(L/L_{\odot}$). \textbf{2 pts; 1 pt for trying; 1 pt for correct answer}


\vspace{8pt}
 \noindent ($iii$) Use the mass-radius and mass-luminosity relations for
 main-sequence stars to express the nuclear timescale of main-sequence stars as
 a function of the mass of the star only. \textbf{2 pts; 1 pt for trying; 1 pt
 for correct answer}

 \vspace{8pt}
 \noindent ($iv$) Describe the meaning of the nuclear timescale in your own
 words. \textbf{2 pts; should be coherent}

 \vspace{10pt}
 \noindent \textbf{Problem \#1b: The thermal timescale ($\tau_{\rm{therm}}$): 6
 points total}

 \vspace{4pt}
 \noindent ($i$) Derive an expression for the thermal timescale in solar units,
 i.e. expressed in terms of $(R/R_{\odot}$), $(M/M_{\odot}$), and
 $(L/L_{\odot}$). \textbf{2 pts; 1 pt for trying; 1 pt for correct answer}

 \vspace{8pt}
 \noindent ($ii$) Use the mass-radius and mass-luminosity relations for
 main-sequence stars to express the thermal timescale of main-sequence stars as
 a function of the mass of the star only. \textbf{2 pts; 1 pt for trying; 1 pt
 for correct answer}

 \vspace{8pt}
 \noindent ($iii$) Describe the meaning of the thermal timescale in your own
 words. \textbf{2 pts; should be coherent}

 \vspace{10pt}
 \noindent \textbf{Problem \#1c: The dynamical timescale ($\tau_{\rm{dyn}}$): 6
 points total}

 \vspace{4pt}
 \noindent ($i$) Derive an expression for the dynamical timescale in solar
 units, i.e. expressed in terms of $(R/R_{\odot}$), $(M/M_{\odot}$), and
 $(L/L_{\odot}$). \textbf{2 pts; 1 pt for trying; 1 pt for correct answer}

 \vspace{8pt}
 \noindent ($ii$) Use the mass-radius and mass-luminosity relations for
 main-sequence stars to express the dynamical timescale of main-sequence stars
 as a function of the mass of the star only. \textbf{2 pts; 1 pt for trying; 1
 pt for correct answer}

 \vspace{4pt}
 \noindent ($iii$) Describe the meaning of the dynamical timescale in your own
 words. \textbf{2 pts; should be coherent}


 \vspace{10pt}
 \noindent \textbf{Problem \#1d: Comparing the timescales: 9 points}

 \vspace{4pt}
 \noindent ($i$) Summarize your results for the questions above by computing the
 nuclear, thermal and dynamical timescales for a $1$, $10$ and $25\,M_{\odot}$
 main-sequence star. \textbf{3 pts; 1.5 pt for trying; 0.5 pt for each correct
 answer}

 \vspace{4pt}
 \noindent ($ii$) When the Sun becomes a red giant, its radius will increase to
 $200\,R_{\odot}$ and its luminosity to $3000\,L_{\odot}$. Estimate
 $\tau_{\rm{dyn}}$ and $\tau_{\rm{therm}}$ for such a red giant. \textbf{4 pts;
 1 pt for trying each and 1 pt for correct}

 \vspace{4pt}
 \noindent ($iii$) How large would such a RG have to become for $\tau_{\rm{dyn}}
 > \tau_{\rm{therm}}$? Assume both R and L increase at constant effective
 temperature. \textbf{2 pts; 1 pt for trying and 1 pt for correct}


 \vspace{10pt}
 \noindent \dunderline{1.0pt}{\textbf{Problem \#2: Accessing scientific
 manuscripts as part of astronomy research: 13 points total}}

 \vspace{4pt}
 \noindent The astrophysics data system
 (\href{https://ui.adsabs.harvard.edu/}{ads}) hosted by the Smithsonian
 Astrophysical Observatory is the central hub of all science manuscripts in
 astronomy. As part of Project 1, you'll likely need to do some searching to
 find science papers that are relevant to your topic. This problem is designed
 to help you get acquainted with ads and to begin your search through the
 astronomical literature!

 \vspace{10pt}
 \noindent \textbf{Problem \#2a: The ads search form; 3 points: search term,
 number, answer to last question}\\
 \noindent The ads search form contains several examples for searching the
 literature, including a full text search. The full text search can be very
 powerful but only if you wield it smartly. Decide on a search term that
 pertains to your project, report it as an answer, then also report the number
 of manuscripts that are returned in your search. How useful is this set of
 literature in bulk?

 \vspace{10pt}
 \noindent \textbf{Problem \#2b: Narrowing down a search; 5 points: 1 per search
 term, 1 per number, 1 for description}\\
 \noindent Often times, the literature is \emph{too much} to sort through. In
 this case, you need to use more specific searches to narrow down the results.
 Based on the available search terms, choose two additional ways to narrow down
 your search. Report those below and also report the number of manuscripts
 returned by the search with each additional narrowed down search. Describe in
 words why you chose the search terms you chose. 

 \vspace{10pt}
 \noindent \textbf{Problem \#2c: Selecting papers to read; 5 pts: 2 per paper
 selection, 1 for answer to question}\\
 \noindent Even with a more manageable set of papers, it can be difficult to
 choose which papers to start reading. One option is to read the abstracts of
 the papers returned starting from the most recent until you find a paper that
 seems relevant. Alternatively, you can start from the oldest paper and read
 abstracts until you find a paper that seems relevant. Try out both of the
 options and pick one paper from each option and report the title and year.
 Which was the easiest option for selection? Why? (Note there isn't a right
 answer but you should be able to articulate which one you prefer since it will
 help you in future literature searches!) 



\end{document} 
